\section{Erotetic search scenarios: complete proofs}
\label{sec:p18-erotetic-search-annex}

Fix the source-defined family conditions and partition semantics of
Section~\ref{sec:p18-erotetic-search-scenarios}. A question implication warrant
always supplies both clauses and an actual earlier question occurrence; it
does not supply unrestricted transitivity. The original continuation states
several results with proof referrals or sketches. The proofs below establish
the declared results directly and keep the specified scope restrictions visible
\cite[pp.~30--43]{Wisniewski2001QuestionsInferences}.

\begin{annexitem}{annex:p18.es.structure}{thm:p18-es-structure}{Prefix structure, first query and finite answer branching}
\begin{proof}
The query definition requires a question. At a query prefix, response fairness
represents every member of its answer set, which contains at least two distinct
formulas. Distinct responses require distinct successor nodes, so this occurrence
branches. At a declarative prefix the second family clause forces one common
successor on every extending path, so no declarative occurrence branches.
The same clause forces one successor at any question followed by another question
or by a nondirect declarative. Such a prefix cannot simultaneously have an
answer successor on another path: the forced-successor clause would be violated.
Hence all genuine forks above the root are query occurrences.

The successor formulas at a query are the image of the finite subfamily of
paths through it. Response fairness says this image includes the entire direct
answer set. Each extending path has exactly one immediate response there, so
the answer set has cardinality at most the number of paths. In particular it
is finite, without an assumption that every auxiliary question in the language
has finitely many direct answers.

No path can be a proper prefix of another. If one ended with principal answer
\(A\) while another continued from the same prefix, that latter path's
declarative-successor clause would require a successor also on the ended path.
This contradicts terminality and the well-typed successor requirement. Two
distinct paths therefore have a first different successor. Since the family
has at least two paths, some fork exists. Before the first fork, all paths
follow the same successors and none can terminate earlier, by prefix-freeness.
Under root separation the root is not an answer fork. The first fork is thus
an auxiliary query shared by every path. No earlier query could exist, since
every query forks. This gives the common first query and major branching point.

Finally, a finite family of finite sequences has finitely many prefixes. Their
lengths are bounded by the maximum path length, giving the claimed finite tree
and height. Root separation, rather than an unproved ban on tautological
principal answers, is what excludes the degenerate first-level answer case.
\end{proof}
\end{annexitem}

\begin{annexitem}{annex:p18.es.definition-boundaries}{thm:p18-es-definition-boundaries}{Literal root and canonical-overlap counterfamilies}
\begin{proof}
Let \(A=p\lor\neg p\), \(B=q\lor\neg q\) for distinct variables,
and \(Q=?\{A,B\}\). Take \(X=\varnothing\) and
\(\Phi=\{(Q,A),(Q,B)\}\). Each final formula is a direct answer and
is entailed by the empty set, whose members all precede it vacuously. Thus
each sequence meets the literal derivation conditions if the consequence
support may be empty. There are no auxiliary questions. The initial-premise
disjointness condition holds, the family has two elements, and at index 1
each direct answer has its stipulated response path. These are the only
proper prefixes. Hence the literal family conditions hold. But no index
satisfies \(1<k<n\), so there is no query. Both paths have a principal
answer in position 2. This verifies the exact additional convention needed
for Facts 6--7; it does not affect Golden Path, since both endpoints are true.

For the canonical-overlap case let \(B=p\lor q\),
\(Q=?\{p,q,r\}\), \(R=?\{B,r\}\), \(X=\{B\}\), and use
\[
 (Q,B,R,B,?p,p),\quad
 (Q,B,R,B,?p,\neg p,q),\quad
 (Q,B,R,r).
\]
The question \(R\) is implied by \(Q\) on \(X\): target coverage holds
because \(B\) is already a premise; witnesses for its two answers are
\(\{p,q\}\) and \(\{r\}\). The later \(?p\) is implied by the
earlier \(Q\) on \(\{B\}\): polar coverage is automatic, positive
response gives witness \(\{p\}\), and negative response plus \(B\)
gives witness \(\{q\}\). All auxiliary answer sets differ from \(dQ\).
The last \(q\) is a consequence of \(B,\neg p\); other terminal
formulas are immediate responses. Every response to \(R\), and then to
\(?p\) on its positive-B branch, is represented. All nonanswer successors
are fixed on shared prefixes. Thus this is a scenario.

Its response \(B\) occurs after the first branch and belongs to \(X\).
Leaving it there violates literal canonical form. Moving it before \(R\)
removes the immediate B-response edge. Paths through the \(R\) prefix
would then have successor \(?p\) on the former B-branch but successor
\(r\) on the other branch. The nondirect-successor clause on the first
path forces \(?p\) on every shared-prefix path, which the r-path violates.
Inserting a duplicate B-premise earlier does not repair that missing edge.
Thus the move-only recipe is not an unrestricted construction. We have not
claimed that a new scenario with different queries cannot exist.
\end{proof}
\end{annexitem}

\begin{annexitem}{annex:p18.es.golden-path}{thm:p18-es-golden-path}{A full generalized Golden Path proof}
\begin{proof}
Fix the admissible partition from the theorem and build a prefix of a path.
Maintain the invariant that every declarative already on the prefix is true
and every question already on it is sound. It holds at the root by assumption.
If the prefix is terminal, its last formula is a direct answer to the principal
question and is true by the invariant; we are done.

Otherwise inspect a path extending the prefix. If its current question is
followed by a direct answer, soundness supplies at least one true direct answer
\(B\). Family fairness gives a path extending exactly the current prefix
with \(B\) as successor. Choose that extension. It preserves the invariant.
This applies also at the literal root-answer edge, because the family fairness
clause includes index 1, regardless of the narrower query convention.

In every other case the family clause fixes a common next expression. If it is
a declarative, its derivation reason is an initial premise, an immediate
auxiliary response, or a consequence of earlier declaratives. The initial
premise is true by the theorem's hypothesis. The response case belongs to the
answer choice already considered. For the consequence case, all its supporting
declaratives are true by the invariant, so semantic entailment makes it true.

If the next expression is a question, its derivation warrant is implication
from an earlier question on a set of earlier declaratives. That earlier
question is sound and every member of the warrant's basis is true. Choose a
true direct answer to that earlier question; the first implication clause
gives a true direct answer to the next question. Hence it is sound. No
composition of ordinary implication has been assumed: each warrant has its
own actual earlier premise and true basis on the constructed prefix.

At each stage the prefix is extended by one expression and remains a prefix
of some scenario path. The finite tree has maximum length \(h_\Phi\), so this
process cannot continue beyond length \(h_\Phi\). It must reach a terminal path.
The maintained invariant proves all three conclusions. Neither completeness
of the endpoint set nor accessibility of answers is used.
\end{proof}
\end{annexitem}

\begin{annexitem}{annex:p18.es.execution}{thm:p18-es-execution}{Exactly what the finite execution bound assumes}
\begin{proof}
Use the same invariant as in Golden Path. A run that receives any true direct
response chooses a represented successor by fairness. Fixed successors have
the initial-premise, semantic-consequence or soundness-transmission justification
proved there. Thus every truthful choice, not only one specially selected
choice, preserves the invariant. Every step increases prefix length by one,
so after at most \(h_\Phi-1\) such steps from the root a terminal direct answer
is reached and is true.

This counts transitions in the given finite object. If a query produces no
response, there is no next transition and the induction does not describe a
completed run. If it produces a false response, the truth invariant need not
hold. The proof contains no algorithm for constructing a finite scenario,
for obtaining true responses, or for bounding the time required by a
declarative consequence step. Those are precisely the separate assumptions
and nonclaims in the corollary.
\end{proof}
\end{annexitem}

\begin{annexitem}{annex:p18.es.basis-replacement}{thm:p18-es-basis-replacement}{Premise substitution without question transitivity}
\begin{proof}
Take a partition satisfying \(Z\) and a source answer \(A\in dQ\).
Every member of \(Y\) is true by the stipulated entailments. The first
clause of \(\operatorname{Im}(Q,Y,R)\) therefore supplies an answer to
\(R\). For any target answer \(B\), retain its original nonempty proper
witness \(W_B\subsetneq dQ\). At a partition satisfying \(Z,B\), all
members of \(Y\) are again true, so some member of \(W_B\) is true.
Thus both clauses hold on \(Z\), with the very same witness sets. This
argument uses truth of the replaced premises, not an implication chain
between new question premises.
\end{proof}
\end{annexitem}

\begin{annexitem}{annex:p18.es.canonical-concise}{thm:p18-es-canonical-concise}{The complete scoped normalization construction}
\begin{proof}
Only finitely many expressions occur in the finite family, so the set \(X_0\)
of its occurring initial premises is finite. Choose one fixed enumeration
without repetitions. On every path insert \(X_0\) immediately after the
root, then delete all original occurrences of formulas in \(X\).
No endpoint is deleted because \(dQ\cap X=\varnothing\). No query response
is deleted because query answer sets avoid \(X\). Moreover an initial-premise
occurrence cannot immediately follow an auxiliary question as a nondirect
declarative: the derivation adjacency clause forbids that. Thus deletion
does not break a query edge or the question-next-question/answer condition.

All former implication bases and declarative consequence supports are still
available earlier: every deleted initial premise is in the inserted common
prefix; other earlier formulas remain in their order. The root-next expression
is not a principal answer, since inserted premises avoid \(dQ\), or the
original root-separated prefix remains when \(X_0\) is empty. Every path
therefore remains a derivation and every occurring formula of \(X\) is now
in the common prefix before the first query.

For fairness, the first point where two original paths differ is a query
response, by the structural proof. That response and all other query
responses are retained, so distinct paths cannot collapse. Moving fixed
initial entries to one common prefix preserves each answer branch and its
full collection of responses. Shared nonquery prefixes still have the same
successor; no new fork is introduced by one identical inserted list. This
proves the scenario family conditions and canonical form.

Now begin with a canonical scenario and delete only intermediate declaratives
outside the three concise classes. Every deleted formula \(D\) is not an
initial premise or immediate response, so its derivation reason is entailment
from earlier declaratives. Induct along each original path: retained earlier
declaratives entail all earlier deleted declaratives, because a partition
making the retained ones true makes the supporting sets true at each induction
step. It therefore makes \(D\) true. This proves the required support
replacement for every deleted occurrence, with no unspecified closure lemma.

For each remaining question, replace its old auxiliary basis by all retained
declaratives preceding it. These entail every deleted member of the old basis;
the basis-replacement lemma supplies its same implication warrant from the
same retained earlier question. Remaining derived endpoints have the analogous
semantic consequence proof. An auxiliary question's immediate response is
never deleted, and an immediate nonquery successor was a question, not one
of the deleted declaratives. All query responses and endpoints are retained.

The first different successors of distinct paths are still retained responses;
no paths collapse and no distinct response histories become identical prefixes.
Hence every response remains represented and deterministic successors remain
deterministic. Initial premises stay in the canonical common prefix. The
result is therefore a concise scenario with the stated endpoints.
\end{proof}
\end{annexitem}

\begin{annexitem}{annex:p18.es.imperative}{thm:p18-es-imperative}{Projection preserves the truthful route, not all question warrants}
\begin{proof}
Golden Path supplies a path of the concise scenario with all declaratives
true, all questions sound and true terminal principal answer. Delete its
unqueried auxiliary question occurrences as prescribed. Every retained
declarative is unchanged and true; every retained question remains sound;
the principal question and endpoint remain. This is a path of the imperative
counterpart, proving its projected Golden Path conclusion. Every branch
response is retained, so the conditional instructions still include all
query responses, not merely the selected truthful path.

Deletion can remove the only earlier question supplying a retained query's
implication warrant. The predicate countermodel proved below supplies a
specific failure of clause (3) despite the correctness of every operational
branch. Thus no claim that the projected sequences are derivations follows
from projection of the truthful path.
\end{proof}
\end{annexitem}

\begin{annexitem}{annex:p18.es.meeting}{thm:p18-es-meeting}{All three meeting branches and exact warrants}
\begin{proof}
The source answers \(\neg p,p\land q,p\land\neg q\) each entail an
answer to \(?p\). Target response \(\neg p\) has singleton witness
\(\{\neg p\}\); response \(p\) has witness
\(\{p\land q,p\land\neg q\}\), since Boolean q-or-not-q coverage
makes one of those true. Both are nonempty proper subsets of the three-answer
set. Hence \(Q\) implies the first query.

On the positive response, use basis \(\{p\}\) and earlier \(Q\) to
warrant \(?q\). Polar target coverage holds. Response \(q\), together
with \(p\), entails the singleton \(\{p\land q\}\); response
\(\neg q\) similarly entails \(\{p\land\neg q\}\). This establishes
the second query warrant without using \(?p\) as a pure implication premise.

The negative first response is already a principal direct answer. The positive
first response and either second response entail the corresponding conjunctive
endpoint. These verify every declarative's derivation reason. Auxiliary
question answer sets have two members, unlike the three principal answers,
and the adjacency clauses hold. The first query has its two response branches;
the second, on the p-prefix, has both of its response branches. All other
successors are fixed. The three paths are distinct, initial premises are
empty, and all three principal answers occur as endpoints. This proves
scenariohood, completeness, canonical form and conciseness.
\end{proof}
\end{annexitem}

\begin{annexitem}{annex:p18.es.destination}{thm:p18-es-destination}{All departure/umbrella branches and completeness distinction}
\begin{proof}
The first target is polar and always sound under Boolean valuation. From
\(X,t\) one obtains \(p\lor q\), the proper witness \(\{p,q\}\)
within the four original answers. From \(X,\neg t\) one obtains \(s\),
the proper singleton witness. Hence \(Q\) implies \(?t\) on \(X\).

On the t-prefix, \(X,t\) already entails \(p\lor q\). Thus every
principal direct answer with that basis covers \(R=?\{p,q\}\); its target
answers p and q have the corresponding singleton witnesses in \(dQ\).
The subquestion \(R\) is therefore legitimately derived and left unasked.
For the next query \(?u\), polar coverage holds. On basis \(X,t\),
\(u\) entails \(q\) by \(q\equiv u\); \(\neg u\) entails
\(\neg q\), and then \(p\) by the already covered \(p\lor q\).
These are singleton proper witnesses in \(dR\), so the earlier \(R\)
warrants \(?u\).

Each endpoint follows from the stated basis and responses: \(s\) from
\(\neg t\supset s,\neg t\); \(q\) from the biconditional and \(u\);
\(p\) from morning coverage, the biconditional and \(\neg u\).
All initial-premise successors are fixed. Both t-responses are represented;
on the positive t-prefix the unqueried R has fixed successor \(?u\), and
both u-responses are represented. The family is a three-path scenario and
its initial premises occur before the first query. Its only other declaratives
are immediate responses and endpoints, so it is concise.

The endpoints are exactly \(s,q,p\), not \(r\), proving incompleteness.
For example, the Boolean assignment \(t=p=r=1\), \(q=u=s=0\) satisfies
\(X\) and makes the omitted answer r true, while the negative umbrella
branch reaches the also-true answer p. Thus the endpoint omission is neither
a logical assertion of \(\neg r\) nor a refutation of Golden Path. An
exclusive-destination interpretation would need separately stated premises.
\end{proof}
\end{annexitem}

\begin{annexitem}{annex:p18.es.binary-warrants}{thm:p18-es-binary-warrants}{The two actual question legs and failure of the shortcut}
\begin{proof}
At every Boolean valuation exactly one of the four displayed conjunctive
alternatives is true, according to the truth values of \(A\) and \(B\).
Thus any answer to the principal polar compound covers \(?\pm|A,B|\).
Each conjunctive target answer fixes the two values and thereby the value
of \(A\otimes B\); it entails that formula or its negation. The resulting
singleton is proper in the principal two-answer set, proving regular
implication to the binary question for each of the four connectives.

For binary question to \(?A\), use witnesses
\(\{A\land B,A\land\neg B\}\) and
\(\{\neg A\land B,\neg A\land\neg B\}\).
Polar target coverage holds, and the Boolean B-or-not-B clause proves each
two-answer witness. They are nonempty proper subsets of the four-answer set.
For \(?B\), use the analogous two subsets indexed by B and not-B.
These are the full second-leg warrants.

To refute the direct shortcut with independent variables, in the conjunction,
implication and equivalence cases fix \(A=1\), and compare \(B=0\) with
\(B=1\). Each compound has opposite values at those two valuations.
Thus target answer A entails neither compound nor its negation, the only
possible singleton proper witnesses in the principal polar answer set.
For disjunction fix \(A=0\) and make the same two B choices: now target
answer not-A entails neither principal alternative. So the direct pure
implication fails in every case, even though both actual binary-question
legs are valid.
\end{proof}
\end{annexitem}

\begin{annexitem}{annex:p18.es.predicate}{thm:p18-es-predicate}{Predicate scenario, all responses and projected-warrant countermodel}
\begin{proof}
At the interpretation of constant a, H gives
\(P_1a\equiv(P_2a\land P_3a)\). Universal instantiation here
means evaluating H at that domain individual; the closed-instance
substitution fact is supplied in Annex~\ref{annex:p18.ei.wh-cases}.
Every model has one of the four binary alternatives. Each alternative with
H entails a principal direct answer: the positive-positive alternative gives
\(P_1a\), and the other three give \(\neg P_1a\). Hence \(Q\) regularly
implies R on \(\{H\}\). The binary-warrant proof gives the pure implications
from R to both actual polar queries. The same earlier R remains present on
the positive P2-prefix, so no direct pure implication from \(?P_2a\) to
\(?P_3a\) is needed.

Response \(\neg P_2a\) and H entail \(\neg P_1a\). Responses
\(P_2a,P_3a\) and H entail \(P_1a\). Response \(\neg P_3a\) and H
entail \(\neg P_1a\). These verify all three terminal consequences.
Both P2-responses are present, and both P3-responses appear after the positive
P2-prefix. H and the unasked binary question have fixed successors; every
auxiliary answer set differs from the principal two-answer set. The family
is therefore a complete scenario, with H before the first query and no
intermediate declarative outside the concise classes.

After deleting R, the first retained query \(?P_2a\) has only Q as a possible
earlier question and only H as an available nonempty declarative basis. In a
one-element model let \(P_2a\) be true and compare two interpretations:
\(P_3a,P_1a\) both true, and \(P_3a,P_1a\) both false. H holds in both.
Target answer \(P_2a\) thus entails neither \(P_1a\) nor \(\neg P_1a\),
so there is no nonempty proper witness in the principal polar answer set.
The same models defeat the empty-basis alternative. Clause (3) of derivation
fails, proving that the imperative counterpart is not the original proof
object. Its operational endpoints remain correct by the three consequences
already proved, and its truthful route by the projection theorem.
\end{proof}
\end{annexitem}

\begin{annexitem}{annex:p18.es.search-procedures}{thm:p18-es-search-procedures}{Every original schematic response branch}
\begin{proof}
For Figures 5--7 the full earlier meeting, departure and predicate proofs
apply after the indicated schematic renaming. Their three terminal cases
are respectively \(\neg A,A\land B,A\land\neg B\);
\(D,B,A\) on \(\neg E\), \(E,F\), \(E,\neg F\);
and \(\neg Pc,Pc,\neg Pc\) on \(\neg Rc\), \(Rc,Tc\),
\(Rc,\neg Tc\). All initial-premise formulas and response cases remain
explicit in the sheet. The hidden two-answer departure premise and the
binary predicate premise are supplied in their underlying scenarios, not
silently removed from their proof justifications.

Figure 8 has paths \((?\neg A,?A,\neg A)\) and
\((?\neg A,?A,A,\neg\neg A)\). Double-negation truth supplies the
positive terminal consequence. Every principal answer covers \(?A\), and
each A-response entails the corresponding singleton principal answer,
giving its regular implication warrant and both response branches.

For each compound Figure 9--12, start its underlying paths with
\((?(A\otimes B),?\pm|A,B|,?A)\). The binary-warrant theorem
establishes both these question moves, and also any later \(?B\).
Append the following complete response/endpoint suffixes:
\[
\begin{array}{c|l}
 \land& (\neg A,\neg(A\land B));\quad
 (A,?B,B,A\land B);\quad(A,?B,\neg B,\neg(A\land B))\\
 \supset& (\neg A,A\supset B);\quad
 (A,?B,B,A\supset B);\quad(A,?B,\neg B,\neg(A\supset B))\\
 \lor& (A,A\lor B);\quad
 (\neg A,?B,B,A\lor B);\quad(\neg A,?B,\neg B,\neg(A\lor B))\\
 \equiv& (A,?B,B,A\equiv B);\quad
 (A,?B,\neg B,\neg(A\equiv B));\\
 & (\neg A,?B,B,\neg(A\equiv B));\quad
 (\neg A,?B,\neg B,A\equiv B).
\end{array}
\]
Each endpoint follows from its displayed response values by the corresponding
Boolean truth clause: conjunction is true only on two positive values;
implication false only on positive/negative; disjunction false only on two
negative values; equivalence true exactly on agreement. This checks every
branch separately, including each short-circuit branch.

All actual polar responses are represented. The binary question is unqueried
and has a fixed successor. Other nonresponse successors and endpoint consequences
are deterministic. Auxiliary answer sets differ from the principal set by
their syntactic forms or cardinalities. Thus these are finite underlying
scenarios on the empty basis; deleting only the binary question gives the
printed imperative procedures. Their three, three, three and four suffixes
are exactly the original response cases, not a truth-table certificate offered
in place of the scenario. Correctness is conditional on true inputs and
responses, as the individual truth-clause proofs show.
\end{proof}
\end{annexitem}
